% ---------------------------------------------------------------------
%  附录 F 参考数据清单
% ---------------------------------------------------------------------
\chapter{参考数据清单}
\label{app:data}

参考数据由 NESSi 2.0.0 生成，生成器位于 \code{.nessi/kbref}，
通过 \code{generate.sh} 调用。文件分两组：\code{data/test} 用于
集成测试，\code{data/tutor} 用于算例回归。

\section{集成测试数据}
\label{sec:data-test}

\begin{table}[H]
    \centering
    \caption[\code{data/test} 的分组]{\code{data/test} 的分组}
    \label{tab:data-test}
    \small
    \begin{tabularx}{\textwidth}{@{}lrlX@{}}
        \toprule
        位置 & 文件数 & 主题 & 内容 \\
        \midrule
        顶层 & 131 & 基础工具与卷积 &
            容器、权重、卷积、微分、变换、稳态与移动窗等 \\
        \code{phase3/} & 32 & Dyson 与平衡态 & 自由传播子和 Dyson 求解器 \\
        \code{phase4/} & 36 & VIE2 与观测量 & VIE2、气泡图与等时观测量 \\
        \code{phase6/} & 18 & 赝粒子 & 赝粒子 Dyson 与 VIE2 等 \\
        \midrule
        合计 & 217 & & \\
        \bottomrule
    \end{tabularx}
\end{table}

每个文件以 \code{kbref 1} 开头，随后是元数据行与列头。
常见的数据类型包括 \texttt{herm\_matrix}、\texttt{scalar}、
\texttt{function}、\texttt{element}、\texttt{labels}、
\texttt{moving}、\texttt{series} 与 \texttt{pseudo} 等。

\section{算例回归数据}
\label{sec:data-tutor}

\begin{table}[H]
    \centering
    \caption[\code{data/tutor} 的文件]{\code{data/tutor} 的文件}
    \label{tab:data-tutor}
    \small
    \begin{tabularx}{\textwidth}{@{}lX@{}}
        \toprule
        文件 & 内容 \\
        \midrule
        \code{hubbard_2b_quick} & Hubbard 链二阶微扰 \\
        \code{hubbard_gw_quick} & Hubbard 链 $GW$ \\
        \code{hubbard_tpp_quick} & Hubbard 链 TPP \\
        \code{hubbard_*_green} & 对应的 Green 函数数据 \\
        \code{fkm_bethe_quench_quick} & Fal'kov--Kimball 淬火 \\
        \code{holstein_bethe_migdal_quick} & Bethe Holstein，Migdal \\
        \code{holstein_bethe_umig_quick} & Bethe Holstein，uMig \\
        \code{holstein_nambu_migdal_quick} & Nambu，Migdal \\
        \code{holstein_nambu_umig_quick} & Nambu，uMig \\
        \code{holstein_impurity_migdal_quick} & 杂质 Holstein，Migdal \\
        \code{holstein_impurity_umig_quick} & 杂质 Holstein，uMig \\
        \code{trunc_bethe} & 截断传播 \\
        \code{steady_bethe_trunc} & 稳态截断 \\
        \code{gw} & $GW$ 参考 \\
        \code{examples_*} & 稳态教程的三个用例 \\
        \bottomrule
    \end{tabularx}
\end{table}

\section{生成与更新}
\label{sec:data-generation}

生成流程：
\begin{enumerate}[label=(\arabic*)]
    \item 在 \code{.nessi} 中编译 NESSi 与 \code{kbref} 生成器；
    \item 运行 \code{generate.sh}，逐个用例写出文本文件；
    \item 把生成的文件放入 \code{data/test} 或 \code{data/tutor}；
    \item 运行 \code{examples/validate}，确认本库结果与文件一致。
\end{enumerate}
生成器的输入参数（网格、模型参数、随机种子）与测试中的构造参数
必须一致，否则比对没有意义。每个文件的头部记录了这些参数，
测试读取时会核对。
